\documentclass[letterpaper,10pt]{article}

\usepackage{latexsym}
\usepackage[empty]{fullpage}
\usepackage{titlesec}
\usepackage{marvosym}
\usepackage[usenames,dvipsnames]{color}
\usepackage{verbatim}
\usepackage{enumitem}
\usepackage[hidelinks]{hyperref}
\usepackage{fancyhdr}
\usepackage[brazilian]{babel}
\usepackage{tabularx}
\input{glyphtounicode}

\usepackage[sfdefault]{roboto}

\pagestyle{fancy}
\fancyhf{}
\fancyfoot{}
\renewcommand{\headrulewidth}{0pt}
\renewcommand{\footrulewidth}{0pt}
\setlength{\footskip}{5pt}

\addtolength{\oddsidemargin}{-0.5in}
\addtolength{\evensidemargin}{-0.5in}
\addtolength{\textwidth}{1in}
\addtolength{\topmargin}{-.5in}
\addtolength{\textheight}{1.0in}

\urlstyle{same}
\raggedbottom
\setlength{\tabcolsep}{0in}

% Configuração para justificar o texto em blocos quadrados e sem recuo
\setlength{\parindent}{0pt}
\setlength{\parskip}{10pt}

\titleformat{\section}{\Large\bfseries\scshape\raggedright}{}{0em}{}[\titlerule]

\pdfgentounicode=1

\newcommand{\resumeItem}[1]{\item\small{#1}}
\newcommand{\resumeSubheading}[4]{
\vspace{-1pt}\item
  \begin{tabular*}{0.97\textwidth}[t]{l@{\extracolsep{\fill}}r}
    \textbf{#1} & #2 \\
    \textit{#3} & \textit{#4} \\
  \end{tabular*}\vspace{-7pt}
}
\renewcommand\labelitemii{$\vcenter{\hbox{\tiny$\bullet$}}$}
\newcommand{\resumeSubHeadingList}{\begin{itemize}[leftmargin=0.15in, label={}]}
\newcommand{\resumeSubHeadingListEnd}{\end{itemize}}

\begin{document}

\begin{center}
    \textbf{\Huge Gabriel Frigo} \\ \vspace{5pt}
    \small +55 11 95570-2000 $|$ 
    \href{mailto:gabriel.frigo4@gmail.com}{gabriel.frigo4@gmail.com} $|$ 
    \href{https://www.linkedin.com/in/gabriel-frigo-b6727b275/}{linkedin.com/in/gabriel-frigo} $|$ 
    \href{https://github.com/GabrielFrigo4}{github.com/GabrielFrigo4}
\end{center}

\vspace{10pt}

\textbf{Assunto: Candidatura à vaga de Estágio em Desenvolvimento}

Prezada Equipe de Recrutamento,

Escrevo para expressar meu forte interesse na vaga de \textbf{Estágio em Desenvolvimento}. Como estudante de \textbf{Ciência da Computação na UFABC} com foco em engenharia de sistemas e alta performance, busco a oportunidade de aplicar meu rigor analítico e competências técnicas em uma organização que valoriza a inovação algorítmica.

Minha trajetória acadêmica e profissional é marcada pela união entre a lógica matemática abstrata e a implementação técnica "perto do metal". Possuo experiência prática no desenvolvimento de \textbf{APIs RESTful utilizando Python 3} para serviços financeiros na \textbf{QI Tech}, onde também atuei com modelagem de dados e conteinerização de ambientes via \textbf{Docker}. Além disso, domino linguagens de baixo nível como \textbf{C, C++, Rust, Zig e Assembly}, conhecimento que aplico tanto na administração de projetos no \textbf{Green Team Hacker Club} quanto no desenvolvimento de software para robótica competitiva na equipe \textbf{Tamandutech}.

Destaquei-me em projetos desafiadores que exigem alta capacidade de otimização, como o desenvolvimento de um motor 2D para o jogo \textbf{SpaceInvaders}, utilizando \textbf{GPGPU} para processamento de física em tempo real com \textbf{OpenCL e OpenGL}. Minha base analítica é reforçada por pesquisas de Iniciação Científica em \textbf{Fluxos em Redes (CNPq)} e \textbf{Teoria de Ramsey (PICME)}, além de uma sólida trajetória em competições, incluindo a \textbf{Medalha de Ouro e a 13ª posição nacional na OBMEP}.

Acredito que minha versatilidade técnica — que transita entre o desenvolvimento backend moderno e a programação de sistemas de alto desempenho — aliada à minha experiência como \textbf{mentor de algoritmos para a OBI}, me qualifica para colaborar de forma efetiva com seus times de engenharia. Possuo proficiência em inglês nível \textbf{B2 Intermediate}, garantindo fluidez para atuar em ambientes técnicos globais.

Estou à disposição para uma entrevista onde poderei detalhar como minhas competências podem gerar valor imediato aos seus projetos.

Atenciosamente,

\textbf{Gabriel Frigo}

\end{document}
